% REP006 探索：超過の三つの指標への分解（親：Q001）。実験前の事前記録の報告
% コンパイル：uplatex を2回 → dvipdfmx（TEXINPUTS で .claude/templates/progress_report// を参照）
\documentclass[twocolumn, 10pt]{ujarticle}
\usepackage{progress_report}

\newcommand{\Nest}{N_{\mathrm{est}}}
\newcommand{\Nreal}{N_{\mathrm{real}}}
\newcommand{\Conv}{\mathrm{Conv}}
\newcommand{\Cost}{\mathrm{Cost}}

\begin{document}

\reportid{REP006}
\questionid{Q001}
\exploration{CPA制約の超過は，実績CPAが決まるまでのどの部分で生じているか}
\reportdate{20261010}
\summary{%
既存8手法のCPA制約の超過が，どの部分で生じているかを測る。
実績CPA÷制約を，買い方$E$・見積もりの偏り$P$・抽選$L$の三つの指標の積に分け，どの指標が超過を作っているかを手法ごとに見る。
実験は未実施である（学習ベース5手法の検証のやり直しが終わるのを待つ）。
}
\beginverification

% ============================================================
\section{目的}\label{sec:e-purpose}
% ============================================================

既存8手法のCPA制約の超過が，実績CPAが決まるまでの，どの部分で生じているかを知る。
REP005は「どの手法で」超過が起きているかを数えた。ここでは，同じ超過を「どの部分で」生じているかに分ける。そのために，超過を三つの指標に分けて測る（\ref{sec:index}）。

% ============================================================
\section{実験}\label{sec:e-exp}
% ============================================================

\subsection{対象のアルゴリズム}\label{sec:e-algo}

8手法は，REP005と同じである（ルールベースのPID・ABid・Online LP，学習ベースのBC・IQL・CQL・BCQ・TD3\_BC）。いずれも入札係数$\alpha$を決めて入札する。
入札や環境は変えず，確定した設定（\texttt{baseline\_params.json}）で動かして観測するだけである。
評価の単位は\textbf{セル}（手法×位置×1日）で，位置は広告主の番号である。48位置×4日の192セルが，1手法の評価になる（表\ref{tab:env}）。

\begin{table}[H]
\centering
\caption{実験の条件}
\label{tab:env}
\small
\begin{tabularx}{\linewidth}{@{}l>{\raggedright\arraybackslash}X@{}}
\toprule
区分 & 内容 \\
\midrule
固定 & 既定の\texttt{test.gin}，背景47社は同梱の戦略，同じ端末（Linux） \\
対象 & 8手法×48位置×エピソード0〜3（学習ベースはseed 1〜3を別々のセルとして集計），864 run \\
\bottomrule
\end{tabularx}
\end{table}

\subsection{評価指標}\label{sec:e-index}

超過がどの部分で生じているかを，三つの指標で測る。まず前提知識を述べ，次に指標の作り方を説明し，最後に指標から取る量を示す。

\subsubsection{前提知識}\label{sec:pre}

AuctionNetは広告入札のシミュレータで，広告主48社のうち1社を評価する手法に置き換える（「プレイヤー」）。残りは同梱の戦略で動く。1日は48の時間帯で，約50万回の広告機会がある。位置ごとに予算とCPA制約$C$（1件の購入に払ってよい額の上限）を持つ。

1つの機会では，入札額の上位3社が枠を得て，1つ下の順位の入札額を払い（第二価格），枠に応じた確率で露出する。支払いは露出した機会だけに生じる。
露出した機会で購入が起きるかは，3段で決まる（表\ref{tab:three}）。入札者が入札に使えるのは①だけである。

\begin{table}[H]
\centering
\caption{購入が決まる3段}
\label{tab:three}
\small
\begin{tabularx}{\linewidth}{@{}>{\raggedright\arraybackslash}p{1.9cm}>{\raggedright\arraybackslash}X>{\raggedright\arraybackslash}p{1.3cm}@{}}
\toprule
段 & 内容 & 知る人 \\
\midrule
①推定価値 & 購入が起きる確率の見積もり & 入札者 \\
②実際の確率 & ①に雑音を加え0〜1に切った確率 & 環境 \\
③抽選 & ②の確率で購入するか決める & 環境 \\
\bottomrule
\end{tabularx}
\end{table}

実績CPAは，1日の支払いの合計÷購入数である。$C$を超えると超過で，成績に$(C/\text{実績CPA})^2$が掛かる。

\subsubsection{超過を三つの指標に分ける}\label{sec:index}

超過の度合いを，$R=$実績CPA$\div C$で表す。$R>1$が超過である。実績CPAは支払いの合計÷購入数なので，$R=\Cost/\Conv/C$と書ける。

\textbf{購入数を3通りに数える}　$R$が1を超えるのは，払った額が多いか，購入が少ないか，その両方である。ただし「購入が少ない」には，性質の違う2つが混ざる。購入が起きにくい機会を買ったことと，起きにくくない機会を買ったのに，たまたま起きなかったことである。前者は入札で直せるが，後者は直せない。この2つを分けるには，購入数を，情報が増える順に3通りに数えればよい。
自分が落札して露出した機会について，①の和$\Nest$（入札者が見積もった購入数），②の和$\Nreal$（環境にとっての期待購入数），③の和$\Conv$（実際の購入数）を取る（表\ref{tab:n}）。$\Nest$と$\Nreal$は，1日が終わったあとに記録から作る量である。

\begin{table}[H]
\centering
\caption{購入数の3通りの数え方}
\label{tab:n}
\small
\begin{tabularx}{\linewidth}{@{}lX>{\raggedright\arraybackslash}p{1.3cm}@{}}
\toprule
記号 & 数え方 & 見方 \\
\midrule
$\Nest$ & ①推定価値の和 & 入札者 \\
$\Nreal$ & ②実際の確率の和（期待購入数） & 環境 \\
$\Conv$ & ③抽選の結果の和 & 結果 \\
\bottomrule
\end{tabularx}
\end{table}

\textbf{隣り合う比を取る}　3つの数の間のずれを1つずつ切り出すために，$\Cost/C$（1件に$C$まで払えるとしたとき買える購入数）を頭に置き，隣り合う比を取る（表\ref{tab:epl}）。
順は，入札者が決める（$E$）→ 環境が隠している（$P$）→ 偶然が決める（$L$），というシミュレータの因果の順になり，各指標に責任の主体が1つ付く。どの指標も，1が「その段のずれが無い」基準で，1を超えるほど$R$を大きくする向きに働く。

\begin{table}[H]
\centering
\caption{三つの指標}
\label{tab:epl}
\small
\begin{tabularx}{\linewidth}{@{}>{\raggedright\arraybackslash}p{1.45cm}>{\raggedright\arraybackslash}p{1.9cm}>{\raggedright\arraybackslash}X@{}}
\toprule
指標 & 式 & 測るもの（動かす主体） \\
\midrule
$E$ 買い方 & $\Cost/\Nest/C$ & 見積もりのもとで，購入1件あたりの支払いが$C$の何倍か（入札者の決定） \\
$P$ 見積もりの偏り & $\Nest/\Nreal$ & 見積もりが，環境の本当の期待の何倍か（環境の雑音） \\
$L$ 抽選 & $\Nreal/\Conv$ & 期待した購入数が，実際の購入数の何倍か（偶然） \\
\bottomrule
\end{tabularx}
\end{table}

\textbf{掛け合わせると元に戻る}　
\[
E\,P\,L=\frac{\Cost}{\Nest C}\cdot\frac{\Nest}{\Nreal}\cdot\frac{\Nreal}{\Conv}=\frac{\Cost}{\Conv\, C}=R
\]
$\Nest$と$\Nreal$が約分されるので，積は定義から必ず$R$に一致する（残差が無い）。したがって，この分け方に漏れもダブりも無い。これは仮説ではなく，測る道具である。

\textbf{超過を4つに分ける}　超過したセルを，どの指標が作ったかで，漏れもダブりも無く4つに分ける（表\ref{tab:four}）。$L$の向きは，区分ごとに別に報告する。

\begin{table}[H]
\centering
\caption{超過したセルの4区分}
\label{tab:four}
\small
\begin{tabularx}{\linewidth}{@{}l>{\raggedright\arraybackslash}p{2.9cm}>{\raggedright\arraybackslash}X@{}}
\toprule
区分 & 条件 & 意味 \\
\midrule
抽選型 & $EP\le 1$ & 期待の上では制約内で，抽選だけで超過 \\
買い方型 & $EP>1,\ P\le1$ & 偏りは寄与せず，買い方で超過 \\
偏り型 & $EP>1,\ E\le1$ & 買い方は制約内で，偏りで超過 \\
両方型 & $E>1,\ P>1$ & 買い方と偏りの両方が寄与 \\
\bottomrule
\end{tabularx}
\end{table}

\textbf{限界}　途中の量に$\Nest$を選んだのはこちらの選択で，別の中間の量を挟めば各指標の大きさは変わる。購入や支出を見て入札を変える手法では$E$と$L$が独立でなく，$\log E$と$\log L$の相関で確かめる。購入0のセルは分解に入らない。

\subsubsection{取る量}\label{sec:e-feat}

\begin{table}[H]
\centering
\caption{指標から取る量}
\label{tab:feat}
\small
\begin{tabularx}{\linewidth}{@{}l>{\raggedright\arraybackslash}X@{}}
\toprule
種類 & 記録する量 \\
\midrule
指標の分布 & $\log E,\log P,\log L$の分位点，各指標が1を超えたセルの割合 \\
超過の区分 & 4区分のセル数と割合，区分ごとの$\log L$の中央値 \\
指標の関係 & $\log E,\log P,\log L$の相関。同じ位置×日での手法間の$\log P$ \\
抽選の確認 & $z=(\Conv-\Nreal)/\sqrt{\sum p(1-p)}$の平均と標準偏差 \\
\bottomrule
\end{tabularx}
\end{table}

現状の記録には$\Nest$と$\Nreal$が無いので，ティックごとに$\Nest$，$\Nreal$，$\sum p(1-p)$を記録するよう\texttt{github/}に足し（結果は変えない），runを流し直す（約5時間の見込み）。学習ベース5手法の検証のやり直し（T004）が終わるまで，流さない。

% ============================================================
\section{結果}\label{sec:e-result}
% ============================================================

実験は未実施である。

\end{document}
